\section{Sharper bounds after deleting vertices}
\label{sec:four-vertex-deletion}

This section uses the degree-sequence theorem from Section~\ref{sec:four-finite} a second time.
It excludes every reduced graph with at most two doubled pairs and at least eleven vertices.
Its finite arithmetic checks do not enumerate graphs.

Retain $\rho$, $a$, $b$, $\kappa$, and $B$ from Section~\ref{sec:four-finite}.
We also use
\[
 a(1)=\sqrt{2/\rho^5},\qquad
 \delta=\rho^{-2}=\frac{\sqrt3}{2},\qquad
 \chi=\frac2{\sqrt3}+\frac38(\rho/6)^{1/4}.
\]
These are values from the same published theorem, or explicit combinations of those values.

\begin{lemma}[An edge bound after deletion]\label{lem:four-deleted-edge}
Let $T$ be a simple graph with degrees in $\{2,3,4\}$.
Assume that no edge joins two degree-two vertices.
Put $Q=\prod_{v\in V(T)}b(\deg_T(v))$.
Then every edge $e$ satisfies $c_e(T)\le\delta Q$.
\end{lemma}

\begin{proof}
Let $d,e\in\{2,3,4\}$ be the endpoint degrees of an edge.
Delete that edge and apply the path theorem in either orientation.
Its cycle count is at most
\[
 \frac{\min\{a(d-1)b(e-1),a(e-1)b(d-1)\}}{b(d)b(e)}Q.
\]
For the possible unordered degree pairs, the following upper coefficients suffice:
\[
\begin{array}{c|ccccc}
(d,e)&(2,3)&(2,4)&(3,3)&(3,4)&(4,4)\\\hline
\text{coefficient}&\rho^{-3}&\rho^{-2}&\rho^{-2}&\kappa/\rho&\kappa.
\end{array}
\]
Each coefficient is at most $\delta$.
The omitted pair $(2,2)$ cannot occur by hypothesis.
\end{proof}

\begin{lemma}[One-vertex deletion]\label{lem:four-one-vertex}
Let $G$ be a connected loopless four-regular graph on $n$ vertices.
Assume that its $k$ doubled pairs form a matching and all other edges are single.
For a vertex $u$, write $N=G-u$.

If $k=0$, choose any $u$.
Then
\[
 c(G)-c(N)\le2\kappa B,
 \qquad \sum_{C\subseteq N}|E(C)|\le(2n-4)\frac\rho6 B.
\]
If $k\ge1$, choose $u$ incident with a doubled pair.
Then
\[
 c(G)-c(N)\le1+\chi B,
 \qquad \sum_{C\subseteq N}|E(C)|\le(2n-4)\frac{\rho^3}{8}B.
\]
In the second case, $N$ has exactly $k-1$ doubled pairs.
\end{lemma}

\begin{proof}
First let $k=0$.
Every cycle through $u$ uses exactly two incident edges.
Lemma~\ref{lem:core-edge-frequency} therefore gives
\[
 c(G)-c(N)=\frac12\sum_{e\ni u}c_e(G)\le2\kappa B.
\]
The graph $N$ has four vertices of degree three and $n-5$ vertices of degree four.
Its degree product satisfies
\[
 Q=b(3)^4b(4)^{n-5}=\frac{\rho^7}{8}B.
\]
Apply Lemma~\ref{lem:four-deleted-edge} to its $2n-4$ edges.
Since $\delta\rho^7/8=\rho/6$, this gives the first length-sum bound.

Now suppose that $u$ is joined twice to $v$.
Write $x,y$ for its two other neighbours.
The matching hypothesis makes $u,v,x,y$ distinct.
The graph $N$ has degree two at $v$, degree three at $x,y$, and degree four at every other vertex.
Subdivide one edge in each of its $k-1$ remaining doubled pairs.
The resulting simple graph $T$ has degree product
\[
 Q=b(3)^2b(4)^{n-4}\rho^k
   =\frac{\rho^5}{8}B.
\]
Its degree-two vertices are $v$ and the new subdivision vertices.
No two such vertices are adjacent.
Indeed, each subdivision vertex lies on a doubled pair disjoint from $v$.

Every original single edge has the same cycle count in $N$ and $T$.
For each doubled pair, use its unsubdivided edge in $T$.
Its parallel mate has the same cycle count in $N$.
Lemma~\ref{lem:four-deleted-edge} therefore bounds every original edge frequency by $\delta Q$.
There are $2n-4$ original edges.
Since $\delta\rho^5/8=\rho^3/8$, their sum gives the second length-sum bound.

A cycle through $u$ is its doubled-edge cycle or uses two distinct neighbours of $u$.
Consequently,
\[
 c(G)-c(N)=1+2P_{vx}(N)+2P_{vy}(N)+P_{xy}(N).
\]
Apply the path theorem to $T$.
Subdivision preserves these terminal paths, so
\[
 P_{vx}(N),P_{vy}(N)
 \le a(2)b(3)^2b(4)^{n-4}\rho^{k-1},
\]
\[
 P_{xy}(N)\le a(3)b(2)b(3)b(4)^{n-4}\rho^{k-1}.
\]
The resulting coefficient of $B$ equals
\[
 \frac{4a(2)b(3)^2+a(3)\rho b(3)}{b(4)^4\rho}
 =\rho^2+\frac{3b(3)}{8b(5)}
 =\frac2{\sqrt3}+\frac38(\rho/6)^{1/4}
 =\chi.
\]
This proves the remaining inequality.
\end{proof}

\begin{theorem}[Small matching cutoff]\label{thm:four-deletion-cutoff}
Under the hypotheses of Lemma~\ref{lem:four-one-vertex}, we have $c(G)<2^n$ whenever
\[
 k=0,\ n\ge12;\qquad\text{or}\qquad
 1\le k\le2,\ n\ge11.
\]
\end{theorem}

\begin{proof}
Theorem~\ref{thm:four-core-cutoff} covers all $n\ge19$.
It remains to check $12\le n\le18$ for $k=0$, and $11\le n\le18$ for $k=1,2$.
Put
\[
 g=\lfloor(n+3)/2\rfloor,\qquad k'=\max\{k-1,0\}.
\]
The elementary proof of Lemma~\ref{lem:core-short-cycles} only uses upper bounds on degrees.
It therefore also applies to $N=G-u$, giving
\[
 c_2(N)=k',\qquad
 c_\ell(N)\le\frac{(6(n-1)+2k')3^{\ell-3}}\ell\quad(\ell\ge3).
\]
Apply the length cutoff to $N$ and then add the bound for cycles through $u$.
Lemma~\ref{lem:four-one-vertex} yields
\[
 \frac{c(G)}{2^n}\le A+T(3/4)^{(n-k)/4},
\]
where
\[
 A=2^{-n}\left({\bf1}_{k\ge1}+(1-2/g)k'
       +\sum_{\ell=3}^{g-1}(1-\ell/g)
          \frac{(6(n-1)+2k')3^{\ell-3}}\ell\right),
\]
and
\[
 T=\begin{cases}
 2\kappa+(2n-4)\rho/(6g),&k=0,\\
 \chi+(2n-4)\rho^3/(8g),&k\ge1.
 \end{cases}
\]

For exact rational verification, use
\[
 \kappa<\frac{845071}{10^6},\quad
 \rho<\frac{1074570}{10^6},\quad
 \rho^3<\frac{1240807}{10^6},\quad
 \chi<\frac{1398653}{10^6}.
\]
The first three comparisons follow from powers sixteen, four, and four.
For the last comparison, use
\[
 \frac1{\sqrt3}<\frac{577351}{10^6},\qquad
 \frac38(\rho/6)^{1/4}<\frac{243951}{10^6}.
\]
Their squares and sixteenth powers are rational.
Substitute these rational upper bounds into $T$ to obtain $\overline T$.
For all twenty-three stated pairs $(n,k)$, exact arithmetic gives
\[
 A<1,\qquad \overline T^4(3/4)^{n-k}<(1-A)^4.
\]
The largest displayed upper ratio is less than $.993$, attained at $(n,k)=(11,2)$.
This proves every remaining case.
\end{proof}



\begin{proposition}[Two deletions at order eleven]\label{prop:four-eleven-simple}
Every simple four-regular graph $G$ on eleven vertices satisfies $c(G)<2^{11}$.
\end{proposition}

\begin{proof}
Choose adjacent vertices $u,v$.
Put $N=G-u$ and $H=G-\{u,v\}$.
The cycles through $u$ contribute at most $2\kappa B$.
In $N$, the vertex $v$ has degree three.
The first case of Lemma~\ref{lem:four-one-vertex} gives $\delta Q_N=(\rho/6)B$.
Thus the cycles through $v$ in $N$ contribute at most
\[
 \frac12\sum_{e\ni v}c_e(N)\le\frac32\delta Q_N=\frac\rho4 B.
\]

Let $r$ be the number of common neighbours of $u,v$ in $G$.
Then $0\le r\le3$.
The graph $H$ has $r$ vertices of degree two, $6-2r$ of degree three, and $3+r$ of degree four.
Its degree product satisfies
\[
 \frac{Q_H}{B}
 =\frac{\rho^r b(3)^{6-2r}}{b(4)^{8-r}}
 =\frac{\rho^{11-r}}{32}
 \le\frac{\rho^{11}}{32}=\frac{\rho^3}{18}.
\]
Allowing the additional degree pair $(2,2)$ in Lemma~\ref{lem:four-deleted-edge} gives
\[
 c_e(H)\le\delta_2Q_H,\qquad
 \delta_2=\frac{a(1)}{\rho^2}=\sqrt{2/\rho^9}.
\]
Indeed, all other degree pairs have coefficient at most $\delta<\delta_2$.
Since $H$ has fifteen edges,
\[
 \sum_{C\subseteq H}|E(C)|\le\frac56\delta_2\rho^3B.
\]

Use cutoff $g=6$ for the cycles of $H$.
This graph is simple, has nine vertices, and has maximum degree at most four.
Hence $c_\ell(H)\le18\cdot3^{\ell-2}/\ell$ for every $\ell\ge3$.
It follows that
\[
 c(G)\le\left(2\kappa+\frac\rho4+
                   \frac5{36}\delta_2\rho^3\right)B
      +\sum_{\ell=3}^{5}(1-\ell/6)
                   \frac{18\cdot3^{\ell-2}}\ell.
\]
The short-cycle sum is rational.
Use the previous rational bounds for $\kappa,\rho,\rho^3$, together with
\[
 \delta_2<\frac{1023197}{10^6}.
\]
The last comparison follows by eighth powers, since $\delta_2^8=16(3/4)^9$.
An exact fourth-power comparison gives $c(G)/2^{11}<.987$.
The certificate script records that comparison without rounded arithmetic.
\end{proof}

The file \texttt{four\_regular\_vertex\_deletion\_certificates.py} verifies all constant comparisons and twenty-four finite inequalities.
Its output is \texttt{four-regular-vertex-deletion-results.json}.
Together with Theorem~\ref{thm:four-deletion-cutoff}, Proposition~\ref{prop:four-eleven-simple} excludes every order at least eleven when $k\le2$.
